% Workaround para criação de capítulo não-numerado alinhado à margem
\chapter*{}
\noindent
\phantomsection{\MakeUppercase{\textbf{Introdução}}}
\addcontentsline{toc}{chapter}{INTRODUÇÃO}

% \newline
% \newline

%\section{Motivação}

\vspace{1.4cm}


A gestão de processos desempenha um papel essencial para a competitividade das empresas no cenário contemporâneo. Mais do que registrar informações, os sistemas digitais tornaram-se o elo que conecta setores internos, clientes e fornecedores, atuando como a espinha dorsal operacional de muitos negócios. Segundo \citeonline{davenport1998}, processos bem-estruturados e apoiados por tecnologia elevam a eficiência organizacional, reduzem erros, diminuem custos e ampliam a capacidade de tomada de decisão. Nesse contexto, os Sistemas Integrados de Gestão Empresarial, ERP (\textit{Enterprise Resource Planning}), consolidam-se como ferramentas indispensáveis para organizações que buscam otimizar fluxos de trabalho, garantir padronização e aumentar sua produtividade.

Entretanto, embora ERPs sejam amplamente difundidos em grandes empresas, muitos setores tradicionais ainda apresentam dificuldades para adotar soluções tecnológicas eficientes, seja por limitações financeiras, seja pela falta de sistemas realmente adaptados às suas realidades. A indústria da marmoraria é um exemplo claro dessa lacuna. Apesar de lidar diariamente com operações complexas, como orçamentos detalhados, medição de peças, controle de produção, logística, estoque de chapas e atendimento ao cliente, grande parte das marmorarias brasileiras ainda depende de processos manuais, planilhas descentralizadas e comunicação informal entre equipes.

Essa fragmentação operacional gera retrabalhos, atrasos, perda de informações e baixa previsibilidade. De acordo com \citeonline{slack2008}, falhas de comunicação são um dos principais fatores responsáveis por desperdícios e gargalos em operações. No caso das marmorarias, isso significa desde erros em medidas e pedidos até falhas de agendamento e desorganização na produção, comprometendo a experiência do cliente e a rentabilidade da empresa.

Ao mesmo tempo, vivemos uma fase em que empresas de todos os portes buscam digitalizar sua administração de forma prática e acessível. A facilidade proporcionada por sistemas em nuvem, interfaces de programação de aplicações, as APIs (\textit{Application Programming Interfaces}), \textit{pipelines} automatizados e arquiteturas escaláveis cria um ambiente muito mais favorável para que pequenos negócios adotem soluções antes inacessíveis. Como defendem \citeonline{pressman2016}, o desenvolvimento de software orientado à eficiência e escalabilidade democratiza o acesso tecnológico e permite transformar rotinas antes limitadas por processos analógicos.

Percebendo esse cenário, este trabalho propõe o desenvolvimento de um ERP voltado especificamente para marmorarias, construído com tecnologias modernas e práticas de engenharia de software atuais. O sistema é oferecido no modelo de \textit{software} como serviço, SaaS (\textit{Software as a Service}): há uma única instalação do sistema, e cada marmoraria que contrata o produto recebe uma conta própria, com seus usuários e dados isolados das demais empresas. Cada empresa contratante é chamada, neste trabalho, de \textit{tenant}, conceito detalhado na Seção~\ref{sec:modelo-multitenant}. O sistema inclui autenticação segura baseada em \textit{JSON Web Token} (JWT), um painel administrativo web, uma página institucional de divulgação (\textit{landing page}) com formulário de cadastro de marmorarias interessadas, gravado no banco de dados PostgreSQL, e um fluxo automatizado de integração contínua e entrega contínua, CI/CD (\textit{Continuous Integration / Continuous Delivery}), executado no GitHub Actions.

A arquitetura foi desenhada para execução em nuvem. As imagens Docker da API são publicadas no GHCR (\textit{GitHub Container Registry}), e o caminho está preparado para implantação em uma máquina virtual na \textit{Google Cloud Platform} (GCP), em um servidor virtual privado, VPS (\textit{Virtual Private Server}), ou em um \textit{cluster} Kubernetes para cenários de maior escala. O armazenamento de mídias da aplicação (logo, avatar e fotos de produtos) foi implementado com um \textit{backend} próprio em disco local, acessado por URLs assinadas, isto é, endereços temporários que carregam uma assinatura criptográfica e só permitem a operação autorizada por um prazo curto. O mesmo mecanismo já está pronto para usar o \textit{Google Cloud Storage} (GCS), serviço de armazenamento de arquivos do Google \cite{gcs-docs}, ou um serviço compatível com a API do Amazon S3 (\textit{Simple Storage Service}) \cite{s3-docs}, como o Cloudflare R2 \cite{r2-docs}, caso a operação em nuvem seja retomada.

Por se tratar de um Produto Mínimo Viável, MVP (\textit{Minimum Viable Product}), de caráter acadêmico, e com o objetivo de manter o material documentado neste trabalho sempre acessível e sob controle de custo durante o período de avaliação, a validação prática foi conduzida em ambiente Docker local, com testes manuais realizados também em dispositivos móveis reais (iPhone e Android). Essa abordagem busca aproximar a realidade operacional das marmorarias da maturidade digital presente em outros setores, entregando um sistema acessível, leve e eficiente.

A solução desenvolvida visa reduzir ruídos de comunicação, centralizar informações, padronizar processos e oferecer uma visão clara da operação para gestores e colaboradores. Ao integrar orçamentos, pedidos, produção e entregas em um fluxo contínuo, o ERP pretende não só otimizar a rotina diária, mas transformar a forma como marmorarias operam, conectando-as a um modelo mais moderno e competitivo.
    
\vspace{0.5cm}
\textbf{Objetivo}
\vspace{0.5cm}

O objetivo deste trabalho é desenvolver um ERP totalmente voltado para as necessidades operacionais de marmorarias, oferecendo uma plataforma centralizada que otimize processos como orçamento, controle de produção, gerenciamento de pedidos, registro de clientes, acompanhamento de serviços e organização do fluxo interno de trabalho. A proposta busca substituir métodos manuais, fragmentados e suscetíveis a erro por uma solução moderna, escalável e segura, capaz de elevar o nível de eficiência e confiabilidade dessas empresas.

Além da construção das funcionalidades essenciais do ERP, este projeto tem como finalidade demonstrar a aplicação prática de tecnologias atuais de desenvolvimento, incluindo autenticação baseada em JWT, empacotamento em contêineres com o Docker, CI/CD com GitHub Actions e publicação de imagens da API no GHCR, com arquitetura pronta para implantação automatizada em uma máquina virtual na GCP, em um VPS, ou em um \textit{cluster} Kubernetes, além de um \textit{backend} de armazenamento de mídias local, com suporte pronto para o GCS ou para um serviço compatível com o Amazon S3, como o Cloudflare R2. Dessa forma, o sistema não apenas atende às demandas do setor, mas também incorpora boas práticas de engenharia de software alinhadas ao estado da arte, reforçando a importância da modernização tecnológica aplicada a pequenos negócios.

A relevância deste trabalho encontra-se na criação de uma solução tecnológica acessível, eficiente e especializada para um nicho que historicamente carece de ferramentas adequadas. Apesar de lidarem com processos extremamente detalhados, como medições precisas, cálculos de cortes, gestão de chapas, etapas de acabamento e prazos de entrega, muitas marmorarias ainda operam com planilhas, anotações avulsas ou sistemas genéricos que não contemplam suas particularidades. Tal realidade gera inconsistências operacionais, desperdício de material, redução da produtividade e aumento de custos, prejudicando a competitividade dessas empresas no mercado. O presente projeto busca, portanto, preencher essa lacuna ao oferecer um sistema realmente alinhado ao fluxo de trabalho do setor.

Do ponto de vista empresarial, um ERP especializado representa uma oportunidade clara de reduzir retrabalhos, padronizar informações, aumentar a previsibilidade da produção e melhorar a relação com clientes. Em um segmento em que erros de medida podem gerar perdas financeiras significativas, a centralização de dados e a automatização de processos tornam-se fatores determinantes para o sucesso. Assim, este trabalho reforça a importância da digitalização como estratégia de crescimento sustentável, permitindo que marmorarias alcancem níveis mais altos de organização e performance.

Sob a perspectiva técnica e acadêmica, o desenvolvimento deste projeto contribui para o entendimento prático de como arquiteturas modernas podem ser aplicadas em sistemas corporativos reais. O uso de \textit{pipelines} automatizados, infraestrutura em nuvem e práticas de segurança atualizadas serve como referência para desenvolvedores e pesquisadores interessados em soluções escaláveis. Além disso, o projeto evidencia a viabilidade de construir sistemas robustos mesmo em contextos de pequenas empresas, democratizando o acesso a tecnologias antes restritas a grandes organizações.

Outro aspecto relevante é a abordagem modular e escalável adotada na construção do sistema. O ERP foi projetado de forma a permitir expansões futuras, como diferenciação de planos por porte da marmoraria (limitando recursos conforme a necessidade), módulos de logística e aprofundamento dos painéis analíticos (\textit{dashboards}). Nesta versão acadêmica, todas as funcionalidades do MVP permanecem disponíveis sem bloqueio por assinatura, priorizando a demonstração completa do fluxo operacional. A arquitetura em camadas facilita futuras melhorias, garantindo que novas funcionalidades possam ser adicionadas sem comprometer a estabilidade do sistema.

Dessa forma, este trabalho não se limita apenas ao desenvolvimento de um software funcional, mas apresenta também uma proposta de transformação operacional para o segmento de marmorarias.

\vspace{0.5cm}
\textbf{Estrutura do Trabalho}
\vspace{0.5cm}

O Capítulo 1 apresenta uma revisão bibliográfica e uma revisão aprofundada da literatura relacionada aos temas centrais do trabalho.

O Capítulo 2 explora, de forma técnica e organizada, todos os conceitos necessários para compreender a construção do sistema.

O Capítulo 3 descreve a metodologia utilizada e detalha o processo de desenvolvimento do ERP.

O Capítulo 4 apresenta os resultados obtidos: testes funcionais, desempenho, qualidade de código e validação da entrega contínua.

O Capítulo 5 discute esses resultados, avalia o alcance dos objetivos definidos nesta introdução e apresenta as limitações identificadas e os trabalhos futuros.

Por fim, a Conclusão e Trabalhos Futuros recapitulam as contribuições do projeto e apontam perspectivas de aprimoramentos.




\vspace{10mm} %5mm vertical space

